\chapter{ทัศนศาสตร์บรรยากาศและทฤษฎีการกระเจิงแสงของอนุภาคฝุ่น PM2.5}
\label{chap:ch02}

\section*{แผนบริหารการสอนประจำบทที่ 2}
\addcontentsline{toc}{section}{แผนบริหารการสอนประจำบทที่ 2}

\noindent\textbf{หัวข้อเนื้อหาประจำบท}
\begin{enumerate}
    \item ทฤษฎีคลื่นแม่เหล็กไฟฟ้าและการลดทอนแสงในบรรยากาศ
    \item การเปรียบเทียบการกระเจิงแบบเรย์ลีห์กับการกระเจิงแบบมี
    \item กฎของคอชมิเดอร์และการสูญเสียระยะการมองเห็น
\end{enumerate}

\noindent\textbf{วัตถุประสงค์เชิงพฤติกรรม}
\begin{enumerate}
    \item อธิบายหลักการ ทฤษฎีฟิสิกส์ และแบบจำลองทางสิ่งแวดล้อมประจำบทที่ 2 ได้อย่างถูกต้อง
    \item วิเคราะห์ สังเคราะห์ และประยุกต์ใช้เครื่องมือตรวจวัดหรือกระบวนการแปรรูปพลังงานได้อย่างมีประสิทธิภาพ
    \item คำนวณ ประเมินผลทางเทอร์โมไดนามิกส์ ทัศนศาสตร์ หรือการลดก๊าซเรือนกระจกได้อย่างถูกต้องตามหลักวิชาการ
\end{enumerate}

\noindent\textbf{การวัดและประเมินผล}
\begin{table}[H]
\centering\small
\begin{tabularx}{\textwidth}{p{0.55\textwidth} p{0.25\textwidth} c}
\toprule
\textbf{เกณฑ์การประเมินผล} & \textbf{เครื่องมือวัดผล} & \textbf{สัดส่วนคะแนน} \\
\midrule
1. ทักษะการปฏิบัติการทดลองและการคำนวณ & แบบประเมินทักษะห้องแล็บ & 40\% \\
2. รายงานผลการทดลองและการวิเคราะห์ & การตรวจดาต้าชีทและรายงาน & 30\% \\
3. แบบฝึกหัดทบทวนและคำถามท้ายบท & แบบทดสอบท้ายบทเรียน & 20\% \\
4. จิตพิสัยและการมีส่วนร่วม & การเข้าเรียนและการตรงต่อเวลา & 10\% \\
\midrule
\textbf{รวมทั้งสิ้น} & & \textbf{100\%} \\
\bottomrule
\end{tabularx}
\end{table}
\newpage

% ==========================================================
% เนื้อหาบทเรียนประจำบทที่ 2
% ==========================================================

\section{ทฤษฎีคลื่นแม่เหล็กไฟฟ้าและการลดทอนแสงในบรรยากาศ}
\index{ทฤษฎีคลื่นแม่เหล็กไฟฟ้าและการลดทอนแสงในบรรยากาศ}

การศึกษาและทำความเข้าใจเกี่ยวกับทฤษฎีคลื่นแม่เหล็กไฟฟ้าและการลดทอนแสงในบรรยากาศ (Atmospheric Extinction and Wave Optics) มีความสำคัญอย่างยิ่งในงานฟิสิกส์สิ่งแวดล้อมและพลังงาน โดยมีรายละเอียดสาระสำคัญดังนี้

เมื่อลำแสงจากดวงอาทิตย์ส่องผ่านชั้นบรรยากาศ ลำแสงจะเกิดการลดทอนความเข้ม (Attenuation หรือ Extinction) อันเนื่องมาจาก 2 ปรากฏการณ์หลัก คือ การกระเจิงแสง (Scattering) และการดูดกลืนแสง (Absorption) โดยโมเลกุลของก๊าซและอนุภาคละอองลอย (Aerosols)

สัมประสิทธิ์การลดทอนแสงรวมของบรรยากาศ (Total Extinction Coefficient, $b_{\text{ext}}$) คำนวณได้จาก:
\[ b\_{\text{ext}} = b\_{\text{scat}} + b\_{\text{abs}} = (b\_{sg} + b\_{sp}) + (b\_{ag} + b\_{ap}) \]
โดยที่ $b_{sg}$ และ $b_{ag}$ คือการกระเจิงและการดูดกลืนโดยก๊าซ, $b_{sp}$ และ $b_{ap}$ คือการกระเจิงและการดูดกลืนโดยอนุภาคฝุ่นละอองตามลำดับ

\begin{physbox}[หลักการและสมการฟิสิกส์สิ่งแวดล้อม]
การทำความเข้าใจพฤติกรรมของอนุภาคและการถ่ายเทความร้อนในระบบสิ่งแวดล้อมต้องอาศัยการประยุกต์ใช้กฎเทอร์โมไดนามิกส์และทัศนศาสตร์คลื่นแม่เหล็กไฟฟ้าอย่างแม่นยำ
\end{physbox}

\section{การเปรียบเทียบการกระเจิงแบบเรย์ลีห์กับการกระเจิงแบบมี}
\index{การเปรียบเทียบการกระเจิงแบบเรย์ลีห์กับการกระเจิงแบบมี}

การศึกษาและทำความเข้าใจเกี่ยวกับการเปรียบเทียบการกระเจิงแบบเรย์ลีห์กับการกระเจิงแบบมี (Rayleigh vs Mie Scattering Regimes) มีความสำคัญอย่างยิ่งในงานฟิสิกส์สิ่งแวดล้อมและพลังงาน โดยมีรายละเอียดสาระสำคัญดังนี้

รูปแบบการกระเจิงของแสงขึ้นอยู่กับ พารามิเตอร์ขนาด (Size parameter, $x$):
\[ x = \frac{2\pi r}{\lambda} \]
โดยที่ $r$ คือรัศมีของอนุภาค และ $\lambda$ คือความยาวคลื่นของแสงที่ตกกระทบ:

1. การกระเจิงแบบเรย์ลีห์ (Rayleigh Scattering; $x \ll 1$ หรือ $r < 0.1\lambda$): เกิดขึ้นเมื่ออนุภาคมีขนาดเล็กกว่าความยาวคลื่นแสงมาก เช่น โมเลกุลของก๊าซไนโตรเจนและออกซิเจน ความเข้มของการกระเจิงแปรผกผันกับความยาวคลื่นยกกำลังสี่ ($I \propto 1/\lambda^4$) แสงสีน้ำเงินจึงกระเจิงได้ดีกว่าแสงสีแดง ทำให้เรามองเห็นท้องฟ้าเป็นสีฟ้า
2. การกระเจิงแบบมี (Mie Scattering; $x \approx 1$ หรือ $0.1\lambda < r < 10\lambda$): เกิดขึ้นเมื่ออนุภาคมีขนาดใกล้เคียงกับความยาวคลื่นแสง (ช่วง 0.1 ถึง 2.5 ไมโครเมตร) ซึ่งตรงกับขนาดของ อนุภาคฝุ่น PM2.5 พอดี การกระเจิงแบบนี้มีความไวต่อความยาวคลื่นน้อยกว่ามาก และมีทิศทางการกระเจิงพุ่งไปข้างหน้า (Forward scattering) เด่นชัด ทำให้ท้องฟ้าในวันที่มีฝุ่นหนาแน่นปรากฏเป็นสีขาวขุ่น มัว และมีทัศนวิสัยต่ำ

\section{กฎของคอชมิเดอร์และการสูญเสียระยะการมองเห็น}
\index{กฎของคอชมิเดอร์และการสูญเสียระยะการมองเห็น}

การศึกษาและทำความเข้าใจเกี่ยวกับกฎของคอชมิเดอร์และการสูญเสียระยะการมองเห็น (Koschmieder's Law of Visual Range) มีความสำคัญอย่างยิ่งในงานฟิสิกส์สิ่งแวดล้อมและพลังงาน โดยมีรายละเอียดสาระสำคัญดังนี้

ความสัมพันธ์ระหว่างความเข้มข้นของฝุ่นละอองในอากาศกับระยะการมองเห็นหรือทัศนวิสัย (Visual Range; $V$) อธิบายได้ด้วย กฎของคอชมิเดอร์ (Koschmieder's Law):
\[ V = \frac{-\ln(C\_{\text{th}})}{b\_{\text{ext}}} = \frac{3.912}{b\_{\text{ext}}} \]
โดยที่ $C_{\text{th}}$ คือเกณฑ์ความแตกต่างของคอนทราสต์ที่ตามนุษย์สามารถแยกแยะวัตถุออกจากฉากหลังได้ ($C_{\text{th}} \approx 0.02$)

เนื่องจากค่าสัมประสิทธิ์การกระเจิงของอนุภาค ($b_{sp}$) มีความสัมพันธ์เชิงเส้นตรงกับความเข้มข้นเชิงมวลของ PM2.5 ตามสมการ $b_{sp} = \alpha \cdot [\text{PM}_{2.5}]$ ดังนั้นระยะทัศนวิสัยจึงแปรผกผันกับความเข้มข้นของฝุ่น PM2.5 โดยตรง ซึ่งเป็นหลักการทางฟิสิกส์พื้นฐานที่นำไปสู่การประเมินค่าฝุ่นจากภาพถ่าย

\section{ขั้นตอนและวิธีปฏิบัติการทดลอง}
\index{ขั้นตอนการทดลองประจำบทที่ 2}

\subsection{การวัดค่าความลึกเชิงแสงของละอองลอย ด้วย Sun Photometer}
\begin{conceptbox}[การวัดค่าความลึกเชิงแสงของละอองลอย ด้วย Sun Photometer]
1. เล็งกล้องตรวจวัดดวงอาทิตย์ (Sun Photometer) ชนิดพกพาเข้าหาศูนย์กลางดวงอาทิตย์
2. วัดค่าความเข้มรังสีตรง (Direct solar irradiance) ที่ความยาวคลื่น 380, 440, 500, 675, และ 870 nm
3. บันทึกมุมเงยของดวงอาทิตย์ (Solar zenith angle) และคำนวณมวลอากาศสัมพัทธ์ ($m$)
4. คำนวณค่า AOD ตามกฎของแบร์-แลมเบิร์ต: $\tau_\lambda = \frac{1}{m} \ln(I_0 / I)$
5. เปรียบเทียบค่า AOD กับความเข้มข้น PM2.5 ภาคพื้นดิน
\end{conceptbox}

\subsection{การสร้างแบบจำลองการสูญเสียคอนทราสต์ของภาพตามระยะทาง}
\begin{conceptbox}[การสร้างแบบจำลองการสูญเสียคอนทราสต์ของภาพตามระยะทาง]
1. ถ่ายภาพทิวทัศน์เมืองที่มีอาคารสูงตั้งอยู่ที่ระยะห่างต่างกัน (500 ม., 1 กม., 3 กม., 5 กม.)
2. คำนวณค่าคอนทราสต์ ($C$) ของอาคารเทียบกับท้องฟ้า: $C = (I_{\text{target}} - I_{\text{sky}}) / I_{\text{sky}}$
3. พล็อตความสัมพันธ์ระหว่างระยะทาง ($d$) กับ $\ln(C)$
4. หาความชันของเส้นตรงเพื่อคำนวณหาค่าสัมประสิทธิ์การลดทอนแสง ($b_{\text{ext}}$)
5. ประมาณค่าทัศนวิสัยสูงสุดของวันตามกฎของคอชมิเดอร์
\end{conceptbox}

\section{ตารางบันทึกผลการทดลองและการอภิปรายผล}
\begin{table}[H]
\centering\small
\caption{ตารางบันทึกผลการทดลองและการตรวจวัดพารามิเตอร์ประจำบทที่ 2}
\label{tab:ch02_datasheet}
\begin{tabularx}{\textwidth}{p{0.3\textwidth} p{0.3\textwidth} X}
\toprule
\textbf{พารามิเตอร์ที่ตรวจวัด} & \textbf{ค่าที่บันทึกได้} & \textbf{การแปลผลและการวิเคราะห์ทางฟิสิกส์} \\
\midrule
การทดสอบชุดที่ 1 & ค่าอยู่ในเกณฑ์มาตรฐาน & สอดคล้องกับแบบจำลองทฤษฎี \\
การทดสอบชุดที่ 2 & ค่ามีแนวโน้มเพิ่มขึ้นตามสัดส่วน & ยืนยันสมมติฐานการวิจัยเชิงประจักษ์ \\
ประสิทธิภาพรวม & ร้อยละ 88.5 & แสดงผลสัมฤทธิ์ในระดับยอมรับได้ \\
\bottomrule
\end{tabularx}
\end{table}

\section{คำถามทบทวนและแบบฝึกหัดท้ายบท}
\begin{enumerate}
    \item จงเปรียบเทียบกลไกการกระเจิงแสงระหว่างโมเลกุลอากาศ (Rayleigh) กับอนุภาค PM2.5 (Mie) ในแง่การพึ่งพาความยาวคลื่นแสง
    \item หากวัดระยะทัศนวิสัยได้ $V = 10$ กิโลเมตร จงคำนวณหาค่าสัมประสิทธิ์การลดทอนแสง $b_{\text{ext}}$
    \item เหตุใดความชื้นสัมพัทธ์ในอากาศ (RH) ที่สูงเกิน 70\% จึงทำให้อนุภาคฝุ่นขยายขนาด (Hygroscopic growth) และส่งผลต่อการกระเจิงแสงอย่างไร
\end{enumerate}
